\begin{annexitem}{annex:p3.adm.12}{thm:p03-adm-satisfaction}{Uniform satisfaction}
\begin{proof}
First consider relational structures. Let \(M\) be the domain of the
coded structure and let \(B=M^{<\omega}\). This is a set in KPU:
finite powers are obtained by recursion on \(\omega\), using products,
and their union is collected by \cref{thm:p03-adm-collection}. Let
\(S\) be the set of subformulas of the given code, obtained by
\cref{thm:p03-adm-syntax}. For each \(\chi\in S\), an admissible
assignment is a member of \(B\) whose length exceeds every index of a
free variable of \(\chi\). If there are no free variables the empty
assignment is admissible. Keeping superfluous coordinates is harmless.

Assign to \(\chi\) the subset \(E_\chi\) of its admissible assignments
at which it is true. For an atom this is bounded Separation from \(B\),
using equality and the coded relations of the structure. Negation takes
the complement relative to the admissible assignments. For a conjunction
of a set \(F\) of formulas, an assignment belongs to \(E_\chi\) exactly
when it belongs, after restriction if necessary, to the truth set of
each constituent. Disjunction uses some constituent. An existential
quantifier tests some element of \(M\) as the value of its bound
variable; a universal quantifier tests every element. One can extend a
finite assignment to the required bound-variable coordinate before the
test. When \(M\) is empty the existential test is false and the universal
test is true, so no extension is needed. When \(M\) is nonempty, unused
coordinates can be filled by the chosen test value. The induction also
shows that the result is independent of these unused coordinates.

All these tests are bounded once \(B\), the constituent family, and
its already computed truth sets are given. Bounded Separation therefore
produces the next truth set. Membership recursion on the tagged formula
code, or its equivalent subformula recursion with the code's rank,
collects the whole evaluation family by
\cref{thm:p03-adm-membership-recursion}. There is no use of the power
set of \(B\): each individual truth set is specified by Separation,
and the family is collected by functional \(\Sigma\)-Replacement.

An evaluation certificate consists of \(S,B\), the evaluation family,
and the certificates for the previously proved syntactic operations.
Its local compositional clauses are bounded. The assertion that such a
certificate puts the given assignment in the root truth set is
\(\Sigma_1\). The same assertion with nonmembership in that set is
also \(\Sigma_1\). Recursion gives existence, and induction on the
code gives uniqueness, so the two assertions are complementary on
legal inputs. This proves \(\Delta_1\) definability. In a transitive
admissible set all bounded checks are absolute, and external induction
on the actual well-founded formula code identifies the result with
Tarskian satisfaction. A nullary atom uses its recorded truth value.
For function signatures one first evaluates finite terms by the same
recursion, or uses the relational graph translation.
\end{proof}
\end{annexitem}

\begin{annexitem}{annex:p3.adm.13}{thm:p03-standard-numbers}{The standard order}
\begin{proof}
In first-order logic with \(Q_0\), require that \(<\) linearly order a
nonempty predicate \(N\), that it have no largest element, and that
\[
 \forall x\bigl(N(x)\to\neg Q_0y\,(N(y)\land y<x)\bigr).
\]
In weak second-order logic replace the last clause by the assertion
that, for every \(x\in N\), there is a finite set whose members are
exactly the predecessors of \(x\) in \(N\). All order clauses are
relativized to \(N\).

For \(x\in N\) let \(r(x)\) be the finite number of its predecessors.
If \(x<y\), those predecessors together with \(x\) are predecessors
of \(y\), so \(r(x)<r(y)\). The finite linearly ordered predecessor
set of an element with \(r(x)=n\) has elements of ranks
\(0,\ldots,n-1\). Thus the range of \(r\) is downward closed in
\(\omega\). It is unbounded because there is no largest element:
otherwise its finite nonempty range, and hence the order itself, would
have a largest element. Consequently \(r\) is an order isomorphism
onto \(\omega\). Define zero as the least element and successor as
immediate succession. Conversely the standard natural order satisfies
the displayed conditions. This proves existence as well as exclusion
of nonstandard orders.
\end{proof}
\end{annexitem}

\begin{annexitem}{annex:p3.adm.14}{thm:p03-recursive-auxiliary-presentation}{Coding an auxiliary model}
\begin{proof}
Let \(S\) name the nonempty distinguished domain. Work in a finite
many-sorted relational signature; it can be written in one sort by
naming disjoint domains and relativizing quantifiers. Keep the original
relations on \(S\). Add a domain \(N\) with the standard-number
presentation of the hypothesis. Add arithmetic relations on \(N\):
addition and multiplication are total functions, relationally coded,
and satisfy their usual zero and successor recursions. External
induction on the now standard \(N\) makes them the actual arithmetic
operations. Hence the ordinary arithmetical definitions of finite
syntax, substitution, variable indices, and computation certificates
have their standard meanings. Only finitely many such definitions are
needed; each is a fixed first-order formula.

A recursively enumerated axiom set has a primitive recursive
certificate relation \(C(e,t)\): \(t\) codes a terminating enumeration
computation that prints the formula with code \(e\). Thus an axiom
code is characterized on \(N\) by \(\exists t\,C(e,t)\). This does
not assume that membership in the enumerated set is decidable.

Add a domain \(B\) of finite assignments, with a length map to \(N\)
and a ternary coordinate relation between \(B,N,S\). A finite list of
first-order clauses requires: every member has a length; its coordinates
at indices smaller than that length exist uniquely and are in \(S\);
it has no other coordinates; an empty list exists; and each list can
be extended by any one member of \(S\), preserving its old coordinates
and appending that member. Since lengths are standard natural numbers,
every code describes an actual finite list. External induction shows
that every actual finite list of elements of \(S\) has a code.
Duplicate codes are permitted. Equality of the coded lists, restriction,
and replacement of a coordinate are first-order definable using the
coordinate relation. In particular every required replacement exists,
since every finite list is represented.

Add a truth relation \(\operatorname{Tr}(e,b)\), with \(e\in N\)
and \(b\in B\). It evaluates only formulas of the original finite
signature, not formulas mentioning \(\operatorname{Tr}\). Use a
fixed syntax with equality, relation atoms, negation, conjunction, and
existential quantification. For a formula code \(e\), let \(m(e)\)
be one more than its largest variable index, or zero when it has none.
Impose the following finitely many universally quantified compositional
clauses whenever the assignment has length at least \(m(e)\):
\begin{enumerate}
\item equality compares the indicated coordinates; for each original
      relation symbol its atomic clause applies that relation to the
      indicated coordinates;
\item truth of a negation is nontruth of its matrix at the same list;
\item truth of a conjunction is truth of both constituents at that list;
\item truth of \(\exists v_i\,\chi\) means that some list of the same
      length, agreeing at every coordinate except possibly \(i\),
      satisfies \(\operatorname{Tr}(\ulcorner\chi\urcorner,b')\).
\end{enumerate}
All subformula indices are bounded by the parent's largest index.
Therefore the clauses apply at every recursive step. A formula code
has a finite standard parse tree. External induction on it proves
that \(\operatorname{Tr}\) is exactly its usual satisfaction relation,
independent of which code represents the same assignment. Nullary
relation clauses require no coordinates. The truth relation on
malformed codes or insufficient assignments is irrelevant.

Finally conjoin
\[
 \forall e\,\forall b\,
 \bigl(\operatorname{Ax}_T(e)\land
       \operatorname{len}(b)\geq m(e)
       \longrightarrow\operatorname{Tr}(e,b)\bigr),
\]
where \(\operatorname{Ax}_T(e)\) abbreviates the computation-certificate
formula above, and all variables are restricted to their named domains.
Every axiom is a sentence. There are assignments of the required length
because \(S\) is nonempty. The established truth equivalence implies
that the distinguished structure satisfies every axiom of \(T\).

Conversely, given a set-sized model of \(T\), adjoin disjoint domains
for \(\omega\) and all its finite assignments, use the actual
arithmetic and coordinate relations, and interpret \(\operatorname{Tr}\)
by the external ordinary satisfaction relation. These are sets in the
surrounding mathematics. They satisfy every clause, including the axiom
clause. The original structure is unchanged. Finitely many original
relation symbols and finitely many syntactic constructors give a finite
signature and a finite conjunction. Combining that conjunction with
the standard-number sentence gives the asserted single sentence.

For KPU the axiom schemas are recursively enumerable in its finite
set-and-urelement signature. Infinity and the assertion of a set of all
urelements are individual first-order axioms. The construction therefore
applies without requiring a truth predicate internal to the KPU model.
The external standard domain does not by itself standardize the model's
internal natural numbers. To impose that additional condition, name
an element \(w\in S\), require in the distinguished structure the
usual first-order assertion that \(w\) is its \(\omega\), and add a
binary relation coding an isomorphism from the ordered domain \(N\)
onto \(\{a\in S:a\mathrel{E}w\}\), ordered by the model's membership
relation \(E\). Being a bijection and preserving and reflecting the
order are finite first-order conditions. They force the internal
natural order to have type \(\omega\). Conversely every model with
standard internal \(\omega\) admits exactly such an external
isomorphism. This proves the final restriction without assuming that
an arbitrary model of KPU is already standard.
\end{proof}
\end{annexitem}
